\section{Safe-by-construction shields}\label{sec:bg-shields}

The canonical formalism is shielding \cite{alshiekh2018shield}: a reactive system synthesised from an LTL safety specification and a finite MDP abstraction by solving a two-player safety game, placed either preemptively or post-posed. It carries proofs my system does not: correctness for all abstraction-consistent environments, minimal interference, and preserved learner convergence. The authors also warn the shield must remain active after learning, which mine does. Two weaker rungs sit below it. Invalid-action masking \cite{huang2020masking} proves the masked update is a valid policy gradient, shows masking scales where penalties fail, and finds agents trained masked remain largely valid unmasked. Neurosymbolic masking \cite{nsam2026} learns the symbolic grounding when constraints are unknown, and its own future work points back towards temporal-logic specifications.

My shield is \emph{not} a shield in the Alshiekh sense: it composes three weaker but sound pieces, enforced by code rather than proved from a specification (\S\ref{sec:des-shield}). There is accordingly no minimal-interference proof and no lookahead: my floor reacts \emph{at} the starvation threshold where a synthesised shield acts before violation becomes unavoidable. The upgrade path is cheap at single-junction scale, and I name it as future work. My topology inverts iLLM-TSC's (\S\ref{sec:bg-llmtsc}, \S\ref{sec:des-overview}), which a 6GB edge budget makes viable; I did not implement the inverse to measure the difference. No prior LLM-based signal-control work states the formal status of its safety layer at all.
